\begin{theorem}[浸入的 Gauss--Codazzi--Ricci 方程]
等距浸入的第二基本形式 $B$、形状算符 $A$ 与法向联络满足 \eqnrefs{eq:math-dg-gauss-general}、\eqnrefs{eq:math-dg-codazzi-general}、\eqnrefs{eq:math-dg-ricci-general}。三式分别来自环境曲率的切向、混合与法向分块。
\end{theorem}

\begin{proof}
先让环境曲率作用在切向量 $Z$。由 Gauss 公式，
\[
\overline\nabla_YZ
=\nabla_YZ+B(Y,Z).
\]
再沿 $X$ 求环境协变导数，并分别对切向与法向部分使用 Gauss--Weingarten 分解：
\[
\begin{aligned}
\overline\nabla_X\overline\nabla_YZ
={}&\nabla_X\nabla_YZ+B(X,\nabla_YZ)\\
&-A_{B(Y,Z)}X
+\nabla_X^\perp B(Y,Z).
\end{aligned}
\]
交换 $X,Y$，再减去 $\overline\nabla_{[X,Y]}Z$。切向部分得到
\[
(\overline R(X,Y)Z)^\top
=R(X,Y)Z
-A_{B(Y,Z)}X+A_{B(X,Z)}Y.
\]
与切向量 $W$ 做内积并用\eqnrefs{eq:math-dg-shape-second-form-relation}，
\[
\begin{aligned}
\overline g(\overline R(X,Y)Z,W)
={}&g(R(X,Y)Z,W)\\
&-\overline g(B(X,W),B(Y,Z))\\
&+\overline g(B(Y,W),B(X,Z)).
\end{aligned}
\]
移项后得到 Gauss 方程
\begin{equation}\label{eq:math-dg-gauss-general}
\boxed{
\begin{aligned}
g(R(X,Y)Z,W)
={}&\overline g(\overline R(X,Y)Z,W)\\
&+\overline g(B(X,W),B(Y,Z))\\
&-\overline g(B(X,Z),B(Y,W)).
\end{aligned}}
\end{equation}

同一个展开的法向部分为
\[
(\overline R(X,Y)Z)^\perp
=(\nabla_X^\perp B)(Y,Z)
-(\nabla_Y^\perp B)(X,Z),
\]
即 Codazzi 方程
\begin{equation}\label{eq:math-dg-codazzi-general}
\boxed{
(\overline R(X,Y)Z)^\perp
=(\nabla_X^\perp B)(Y,Z)
-(\nabla_Y^\perp B)(X,Z).}
\end{equation}

最后让环境曲率作用在法向场 $\xi$。由 Weingarten 公式，
\[
\overline\nabla_Y\xi
=-A_\xi Y+\nabla_Y^\perp\xi.
\]
再次沿 $X$ 求导并取法向部分：
\[
\begin{aligned}
(\overline\nabla_X\overline\nabla_Y\xi)^\perp
=-B(X,A_\xi Y)
+\nabla_X^\perp\nabla_Y^\perp\xi.
\end{aligned}
\]
交换 $X,Y$ 并减去 $[X,Y]$ 项，得到
\[
(\overline R(X,Y)\xi)^\perp
=R^\perp(X,Y)\xi
-B(X,A_\xi Y)+B(Y,A_\xi X).
\]
与另一法向场 $\eta$ 做内积，利用 形状算符 自伴性，
\[
\begin{aligned}
\overline g(\overline R(X,Y)\xi,\eta)
={}&\overline g(R^\perp(X,Y)\xi,\eta)\\
&-g(A_\eta X,A_\xi Y)
+g(A_\eta Y,A_\xi X)\\
={}&\overline g(R^\perp(X,Y)\xi,\eta)
-g([A_\xi,A_\eta]X,Y).
\end{aligned}
\]
因此 Ricci 方程为
\begin{equation}\label{eq:math-dg-ricci-general}
\boxed{
\overline g(R^\perp(X,Y)\xi,\eta)
=\overline g(\overline R(X,Y)\xi,\eta)
+g([A_\xi,A_\eta]X,Y).}
\end{equation}
交换子的次序已经由上述展开固定；若把 $[A_\xi,A_\eta]$ 反写，右端必须同时变号。

退化自检：一维法丛（如 $\mathbb R^3$ 中曲面）中 $[A_\xi,A_\eta]=0$，Ricci 方程 的法向部分自动满足；$\RR^4$ 中的标准 Clifford 环面具有可同时对角化的两个法向形状算符，故它们的交换子为零，法丛平坦；这由具体形状算符检验，而不是把它们笼统记作 $\pm J$。
\end{proof}

\begin{theorem}[Euclidean 浸入的局部基本定理]
在可缩坐标邻域上给定度量 $g$、秩 $q$ 的带度量法丛及度量联络、对称 $B$。若满足平直环境的 Gauss--Codazzi--Ricci 方程，则存在等距浸入 $F:M\to\RR^{m+q}$ 实现这些数据，唯一性差一个 Euclidean 刚体运动。
\end{theorem}

\begin{proof}
在局部正交标架下，把 Levi--Civita 联络、正规 联络 与第二基本形式组合成一个 $\mathfrak{so}(m+q)$-值一形式
\[
\mathcal A=
\begin{pmatrix}
\omega^i{}_j & -h^i{}_{\beta k}\theta^k\\
 h^\alpha{}_{jk}\theta^k & \omega^\alpha{}_{\beta}
\end{pmatrix}.
\]
这里上下混合块用环境度量识别，符号选择与\eqnrefs{eq:math-dg-moving-frame-second-form}一致。计算其曲率
\[
\mathcal F
:=\dd\mathcal A+\mathcal A\wedge\mathcal A.
\]
切--切块恰是 Gauss 方程的左减右，切--法块恰是 Codazzi 方程，法--法块恰是 Ricci 方程。因此
\begin{equation}
\text{Gauss+Codazzi+Ricci}
\quad\Longleftrightarrow\quad
\mathcal F=0.
\label{eq:math-dg-submanifold-flat-extended-connection}
\end{equation}
也就是说，全部兼容条件等价于这个扩展联络平坦。

在单连通区域上，平坦矩阵联络可积分为一个正交矩阵场 $Q:M\to O(m+q)$，使
\[
Q^{-1}\dd Q=\mathcal A
\]
（若采用相反的矩阵 约定，等价地写成 $\dd Q=Q\mathcal A$）。令 $Q_1,\ldots,Q_m$ 为前 $m$ 个列向量，并定义 $\RR^{m+q}$-值一形式
\[
\alpha:=\sum_{i=1}^m Q_i\,\theta^i.
\]
计算外微分：
\[
\dd\alpha
=\sum_i \dd Q_i\wedge\theta^i
+\sum_i Q_i\,\dd\theta^i.
\]
$\dd Q_i$ 的切向联络部分与无挠第一结构方程
\[
\dd\theta^i+\omega^i{}_j\wedge\theta^j=0
\]
恰好抵消；法向部分含
\[
h^\alpha{}_{ij}\theta^j\wedge\theta^i,
\]
又因 $h^\alpha{}_{ij}=h^\alpha{}_{ji}$ 对称而为零。故
\[
\dd\alpha=0.
\]
单连通性使闭的一形式 $\alpha$ 可全局积分，所以存在 $F:M\to\RR^{m+q}$ 满足
\[
\dd F=\alpha.
\]
因为 $Q_i$ 正交归一，
\[
\langle \dd F(X),\dd F(Y)\rangle
=\sum_i\theta^i(X)\theta^i(Y)=g(X,Y),
\]
故 $F$ 是等距浸入。再由 $\dd Q=Q\mathcal A$ 的混合块读出其第二基本形式与 正规 联络 正是预先给定的 $B$ 与 $\nabla^\perp$。这完成了存在性重构的核心计算。

若两个重构具有相同初始点与初始适配标架，它们满足同一一阶系统，因此由唯一性完全一致；改变初始点和初始正交标架只对应环境中的平移与正交变换。于是唯一性恰好差一个 Euclidean 刚体运动。
\end{proof}

\begin{proposition}[平行曲面的平均曲率响应]
对 Euclidean 定向超曲面，取 $F_s=F+sn$、$A=-\dd n$，并限制在 $1-s\kappa_i>0$ 的范围。则\[\operatorname{vol}_{g_s}=\det(I-sA)\operatorname{vol}_g,\qquad\left.\partial_s\operatorname{vol}_{g_s}\right|_{s=0}=-m\mathsf H\operatorname{vol}_g.\]
\end{proposition}

\begin{proof}
由\eqnrefs{eq:math-dg-parallel-area-jacobian}，平行曲面的面积元为
\begin{equation*}
\operatorname{vol}_{g_s}=|\det(I-sA)|\,\operatorname{vol}_g
=\prod_{i=1}^m|1-s\kappa_i|\,\operatorname{vol}_g,
\end{equation*}
其中 $\kappa_1,\ldots,\kappa_m$ 是 形状算符 $A$ 的本征值即主曲率。在 $s=0$ 附近假设尚未遇到 focal point，即 $1-s\kappa_i>0$ 对所有 $i$ 成立，可去掉绝对值，
\begin{equation*}
\operatorname{vol}_{g_s}=\det(I-sA)\,\operatorname{vol}_g
=\prod_{i=1}^m(1-s\kappa_i)\,\operatorname{vol}_g.
\end{equation*}
把行列式按特征多项式在 $s=0$ 处 Taylor 展开。由 $\det(I-sA)$ 是 $s$ 的 $m$ 次多项式，常数项为 $\det I=1$，一次项为 $-s\operatorname{tr}A$（特征多项式系数与初等对称多项式的关系，一次项系数恰为负迹），故
\begin{equation*}
\det(I-sA)=1-s\operatorname{tr}A+O(s^2).
\end{equation*}
等价地，对乘积形式逐项展开，
\begin{equation*}
\prod_{i=1}^m(1-s\kappa_i)
=1-s\sum_{i=1}^m\kappa_i+O(s^2)
=1-s\operatorname{tr}A+O(s^2),
\end{equation*}
其中用了 $\operatorname{tr}A=\sum_i\kappa_i$。由平均曲率定义 $\mathsf H=\frac1m\operatorname{tr}A=\frac1m\sum_i\kappa_i$，
\begin{equation*}
\operatorname{tr}A=m\mathsf H
\end{equation*}
（余维一时 $\mathbf H=\mathsf H n$）。代回面积元展开，
\begin{equation*}
\operatorname{vol}_{g_s}=(1-sm\mathsf H+O(s^2))\,\operatorname{vol}_g.
\end{equation*}
对 $s$ 求导并在 $s=0$ 取值，
\begin{align*}
\left.\frac{\dd}{\dd s}\operatorname{vol}_{g_s}\right|_{s=0}
&=\left.\frac{\dd}{\dd s}\bigl[(1-sm\mathsf H+O(s^2))\,\operatorname{vol}_g\bigr]\right|_{s=0}\\
&=\bigl[-m\mathsf H+O(s)\bigr]\Big|_{s=0}\operatorname{vol}_g\\
&=-m\mathsf H\,\operatorname{vol}_g.
\end{align*}
这正是后面\eqnrefs{eq:math-dg-first-variation-submanifold-volume}在常法向速度 $V=n$ 情形的局部版本。换言之，平均曲率既是 形状算符 的迹，也是平行法向推进时面积元的一阶响应。
\end{proof}

\begin{proposition}[浸入体积的一阶变分]
设 $F_t$ 为浸入变分、$V=\partial_tF$、$\mathbf H=m^{-1}\operatorname{tr}_gB$。则\[\partial_t\operatorname{vol}_{g_t}=\bigl(\operatorname{div}_MV^\top-m\overline g(\mathbf H,V^\perp)\bigr)\operatorname{vol}_g.\]闭子流形或紧支法向变分下，体积临界条件恰为 $\mathbf H=0$。
\end{proposition}

\begin{proof}
取局部坐标 $x^i$，诱导度量为
\[
g_{ij}(t)
=\overline g(\partial_iF_t,\partial_jF_t).
\]
时间求导并用环境联络度量相容性：
\[
\dot g_{ij}
=\overline g(\overline\nabla_iV,e_j)
+\overline g(e_i,\overline\nabla_jV).
\]
因此
\[
\frac12\operatorname{tr}_g\dot g
=g^{ij}\overline g(\overline\nabla_iV,e_j).
\]
将 $V=V^\top+V^\perp$ 分解。切向部分给出
\[
g^{ij}g(\nabla_iV^\top,e_j)
=\operatorname{div}_M V^\top.
\]
对法向部分，利用
\[
0=e_i\,\overline g(V^\perp,e_j)
=\overline g(\overline\nabla_iV^\perp,e_j)
+\overline g(V^\perp,\overline\nabla_ie_j)
\]
得到
\[
\overline g(\overline\nabla_iV^\perp,e_j)
=-\overline g(V^\perp,B_{ij}).
\]
收缩后
\[
\frac12\operatorname{tr}_g\dot g
=\operatorname{div}_M V^\top
-m\overline g(\mathbf H,V^\perp).
\]
再用体积形式的一般变分公式
\[
\partial_t\operatorname{vol}_g
=\frac12\operatorname{tr}_g\dot g\,\operatorname{vol}_g
\]
并对散度项积分，即得到\eqnrefs{eq:math-dg-first-variation-submanifold-volume}。

一阶变分 $\delta\operatorname{Area}=-\int_M m\,\overline g(\mathbf H,V^\perp)\dd\operatorname{vol}_g$ 说明面积泛函的 Euler--Lagrange 方程为 $\mathbf H=0$，即极小子流形方程：切向变分只产生散度（边界项），法向变分才通过平均曲率与面积变化耦合。

典型例：平面 $B_{ij}=0$、$\mathbf H=0$；半径 $R$ 的球面取外法向时 $\mathbf H=-(1/R)n\neq0$；悬链面 $\mathbf H=0$，即两圆环间肥皂膜的平衡形状。
\end{proof}

\begin{theorem}[Euclidean 极小超曲面的 Simons 恒等式]
设 $M^m\subset\RR^{m+1}$ 为光滑极小超曲面，$A$ 为形状算符。则\[\Delta_{\mathrm{LB}}h_{ij}=-|A|^2h_{ij},\qquad\tfrac12\Delta_{\mathrm{LB}}|A|^2=|\nabla A|^2-|A|^4.\]
\end{theorem}

\begin{proof}
在 Euclidean 超曲面 上，Codazzi 给出
\[
\nabla_k h_{ij}=\nabla_i h_{kj}.
\]
因此
\[
\Delta_{\mathrm{LB}}h_{ij}
=\nabla^k\nabla_kh_{ij}
=\nabla^k\nabla_i h_{kj}.
\]
交换两次协变导数时产生 Riemann 曲率 对两个切指标的作用；再用 Gauss 方程把所有内在 Riemann 张量 写成 $h*h$。另一方面，极小 条件
\[
\operatorname{tr}A=0
\]
使 $\nabla_i\nabla_j(\operatorname{tr}A)$ 项完全消失。剩余的所有三次项整理后恰为
\[
-|A|^2h_{ij},
\]
从而得到\eqnrefs{eq:math-dg-simons-hessian}。最后对
\[
|A|^2=h_{ij}h^{ij}
\]
应用乘积法则：
\[
\frac12\Delta_{\mathrm{LB}}|A|^2
=|\nabla A|^2
+\langle\Delta_{\mathrm{LB}}A,A\rangle
=|\nabla A|^2-|A|^4.
\]
所以 Simons 恒等式 的三项来源分别非常清楚：Codazzi 负责交换一个导数，Gauss 负责把交换子曲率变成 $A^2$，minimality 负责杀掉 平均-曲率 Hessian。

交换子计算可在一点的主曲率正交标架中完整核对。
记 $H_0=\operatorname{tr}A$，在该点 $h_{ij}=\kappa_i\delta_{ij}$。
由 Codazzi 与导数交换，
\[
\Delta_{\mathrm{LB}}h_{ij}
=\nabla_i\nabla_jH_0+
\sum_k([\nabla_k,\nabla_i]h)_{kj}.
\]
两个协变指标各贡献一项；由 Gauss 方程，它们分别给出
\[
\sum_k([\nabla_k,\nabla_i]h)_{kj}
=H_0(A^2)_{ij}-|A|^2h_{ij}.
\]
对角分量的逐项计算尤其直接：$K(e_i,e_k)=\kappa_i\kappa_k$，因此
\[
\begin{aligned}
\sum_k([\nabla_k,\nabla_i]h)_{ki}
&=\sum_kK(e_i,e_k)(\kappa_i-\kappa_k)\\
&=\kappa_i^2\sum_k\kappa_k-\kappa_i\sum_k\kappa_k^2\\
&=H_0\kappa_i^2-|A|^2\kappa_i.
\end{aligned}
\]
非对角项由同一 Gauss 分量计算为零；张量性使该表达与标架无关。
所以一般 Euclidean 超曲面满足
\[
\Delta_{\mathrm{LB}}h_{ij}
=\nabla_i\nabla_jH_0+H_0(A^2)_{ij}-|A|^2h_{ij}.
\]
极小条件 $H_0=0$ 才消去前两项。
这里 $\operatorname{Ric}=H_0h-A^2$，且一般 $A^2\ne|A|^2I$；
例如主曲率为 $1,-1$ 时，$A^2=I$ 而 $|A|^2=2$。
因此不能用这个错误的矩阵恒等式来压缩三次曲率项。
\end{proof}
